\section{Projektplanung- und Management}\label{sec:section-one}

In diesem Kapitel wird zu Beginn der bisherige Stand aus den ersten beiden Praxismodulen zusammengefasst.
Der zweite Abschnitt behandelt die Zielsetzung dieses dritten Praxismoduls, gefolgt von der aktuellen Projektorganisation.
Abschließend wird ein Überblick über die verwendeten Tools in diesem Praxismodul gewehrt sowie auf die neue gezielte Arbeitsweise mit verschiedenen git-branches eingegangen.

\subsection{Rückblick auf Praxismodul I \& 2}\label{subsec:subsection-one-one}

Lorenz ipsum dolor sit amet, consetetur sadipscing elitr, sed diam nonumy eirmod tempor invidunt ut labore et dolore magna aliquyam erat, sed diam voluptua.
At vero eos et accusam et justo duo dolores et ea rebum.
Stet clita kasd gubergren, no sea takimata sanctus est Lorenz ipsum dolor sit amet.

\newpage

\subsection{Zielsetzungen im Praxismodul III}\label{subsec:subsection-one-two}

Lorenz ipsum dolor sit amet, consetetur sadipscing elitr, sed diam nonumy eirmod tempor invidunt ut labore et dolore magna aliquyam erat, sed diam voluptua.
At vero eos et accusam et justo duo dolores et ea rebum.
Stet clita kasd gubergren, no sea takimata sanctus est Lorenz ipsum dolor sit amet.


\subsection{Projektorganisationen}\label{subsec:subsection-one-three}

Das Team besteht unverändert aus den vier Studierenden Anika, Darko, Lorenz und Robin.
Projektbetreuer ist dankenswerterweise weiterhin Herr Prof.\, Dr.\, Dirk Schweim.
Die Rollen- und Aufgabenverteilung innerhalb des Teams hat sich nach den Erfahrungen und Aufbau von entsprechendem Know-How aus den beiden vorherigen Praxismodulen inzwischen gefestigt.
Die Organisation und Übersicht der einzelnen Aufgaben wird weiterhin mit einem Projektmanagementprogramm gesteuert, wöchentliche Meetings dienen der Abstimmung und dem Austausch.

Anika und Darko sind für die grafische und technische Frontendentwicklung zuständig, wobei Anika sich auf die grafische Entwicklung basierend auf dem Design konzentriert und Darko die technische Umsetzung in JavaScript übernimmt.
Für die Backend-Entwicklung ist Lorenz zuständig, durch umfangreiche Berührungspunkte mit den Datenbanken übernimmt er ebenfalls deren funktionelle Integration mit dem System.
Zudem ist er direkter Ansprechpartner für Konflikte und Fragen bezüglich git-Branches und der damit verbundenen Arbeitsweise.
Durch seine intensive Einarbeitung in Docker ist Robin für die Containerisierung mit Docker zuständig, zudem ist er für den virtuellen Server und das Deployment der Plattform auf diesem zuständig.
Des Weiteren übernimmt er den Lead für die Dokumentation sowie Präsentation.

\subsection{Verwendete Tools}\label{subsec:subsection-one-four}

Bis auf eine Ausnahme sind die im Projekt verwendeten Tools unverändert geblieben.
Es hat sich gezeigt, dass die im Praxismodul I und II genutzten Tools für die Projektarbeit gut geeignet sind und keine größeren Änderungen notwendig sind.
Der Abschnitt \textit{Verwendete Tools} aus dem Praxismodul I und II wird daher nicht erneut vollständig wiederholt, es wird lediglich auf die Änderungen eingegangen.

Es kam ein neuer Docker-Container hinzu, der als Mailserver fungiert.
Postfix wurde als Mail Transport Agent (MTA) ausgewählt, um E-Mails zu versenden.
Beabsichtigt wird dieser, um der Plattform eine E-Mail-Funktionalität zu ermöglichen, insbesondere für die Account-Verifizierung und Passwort-Wiederherstellung.

Nachfolgend wird eine Übersicht der verwendeten Tools als Grafik gegeben, die im Praxismodul III zum Einsatz kommen.

\begin{figure}[H]
    \centering
    \includegraphics[width=1.0\textwidth]{tools}
    \caption{Übersicht der Tools}\label{fig:uebersicht-tools}
\end{figure}

\subsection{Arbeit mit Pull Requests und Restrictions in Git}\label{subsec:arbeit-mit-pull-requests-und-restrictions-in-git}

Im aktuellen Praxismodul III wurde die bereits im vorherigen Semester eingeführte Arbeitsweise mit separaten Feature-Branches weiterentwickelt und professionalisiert.
Während die Nutzung von Branches bereits in Praxismodul II etabliert wurde, lag der Fokus nun darauf, den Entwicklungsprozess durch verbindliche Strukturen und Regeln qualitativ weiter zu verbessern.

Neu eingeführt wurden sogenannte Pull Requests (PRs), über die Änderungen aus einem Feature-Branch nicht mehr direkt, sondern nur nach expliziter Prüfung durch ein Review in den `main`-Branch übernommen werden dürfen.
Dazu wurden feste Reviewer pro Pull Request definiert, die jeweils für die Freigabe verantwortlich sind.
Bei größeren Änderungen wurde gemeinsame Review Sitzungen während Team Meetings durchgeführt.
Zusätzlich wurden Git-Restrictions aktiviert, welche das direkte Pushen in den `main`-Branch unterbinden.
Dadurch ist sichergestellt, dass keine unkontrollierten Änderungen in die stabile Hauptversion gelangen.

Diese Weiterentwicklung der Git-Nutzung brachte folgende Vorteile mit sich:
\begin{itemize}
    \item Durch Reviews und Restrictions wurden potenziell fehlerhafte oder ungetestete Änderungen frühzeitig erkannt.
    \item Gemeinsame Review Sessions haben zu umfangreicherem Wissenstransfer innerhalb des Teams gesorgt.
    \item Die stabile Version im `main`-Branch konnte nicht durch versehen und unsauberes Arbeiten überschrieben werden.
\end{itemize}

Gleichzeitig entstanden aber auch folgende Herausforderungen:
\begin{itemize}
    \item Im Team musste vorerst ein Verständnis über den neuen Workflow geschaffen werden.
    \item Auch für kleinere Änderungen waren nun zusätzliche Review-Schritte notwendig.
\end{itemize}

Allgemein hat die Einführung der neuen Herangehensweise zu einer strukturierten Arbeitsweise geführt.
